% !TeX root = part8.tex
\documentclass[../final.tex]{subfiles}

\begin{document}

% Конспект: Desktop/6.jpg--8.jpg (24.10). Кривые на рисунке посчитаны
% численно по формулам ниже. Добавлены высокотемпературная поправка
% для цепочки и сравнение показателей степени.
\practice{8}{24.10.2026}{Теплоёмкость кристаллической решётки}
\captionsetup{font=footnotesize,labelfont=bf,skip=5pt}

\begingroup
\renewcommand{\thesection}{4.\arabic{section}}
\task[9]{Теплоёмкость одномерной и трёхмерной решёток}
\label{task:fksv_heat_capacity}

\condition
Определить теплоёмкость одно-, дву- и трёхмерной решёток.

\solution
Полагаем \rassumption{$k_{\mathrm B}=1$}. Нулевая энергия от
температуры не зависит, поэтому учитываем только тепловую часть
\[
    E=\int\frac{\hbar\omega\,g(\omega)\,d\omega}{e^{\hbar\omega/T}-1}.
\]

\subsection*{Одномерная цепочка}

С плотностью состояний прошлого занятия
\[
    E=\frac{2L\hbar}{\pi a\omega_0}\int_0^{\omega_0}
       \frac{\omega\,d\omega}
            {(e^{\hbar\omega/T}-1)\sqrt{1-\omega^2/\omega_0^2}}.
\]

\subsubsection*{Низкие температуры, $\hbar\omega_0\gg T$}
Существенны $\omega\sim T/\hbar\ll\omega_0$, поэтому
$1/\sqrt{1-\omega^2/\omega_0^2}\simeq1+\omega^2/(2\omega_0^2)\simeq1$,
а верхний предел заменяется бесконечностью:
\[
    E\simeq\frac{2L\hbar}{\pi a\omega_0}
       \left(\frac T\hbar\right)^2\int_0^\infty\frac{z\,dz}{e^z-1}
     =\frac{2L\hbar}{\pi a\omega_0}\left(\frac T\hbar\right)^2
       \frac{\pi^2}6.
\]
Так как $L/a=N$,
\[
    C=\frac{dE}{dT}=\frac{2\pi}{3}\,\frac{NT}{\hbar\omega_0}.
\]
Поправка от следующего члена разложения корня имеет порядок
$N(T/\hbar\omega_0)^3$.

\subsubsection*{Высокие температуры, $\hbar\omega_0\ll T$}
Теперь $z\le z_0\ll1$ и $e^z-1\simeq z$:
\[
    E\simeq\frac{2LT}{\pi a}\int_0^{z_0}\frac{dz}{\sqrt{z_0^2-z^2}}
     =\frac{2NT}{\pi}\arcsin\frac z{z_0}\bigg|_0^{z_0}=NT,
    \qquad C=N.
\]
Каждая из $N$ мод даёт классический вклад $T$.
Удерживая $1/(e^z-1)\simeq1/z-1/2+z/12$, получаем
\[
    C=N\left[1-\frac1{24}\left(\frac{\hbar\omega_0}T\right)^2+\ldots\right].
\]

\subsection*{Трёхмерный кристалл: модель Дебая}

Для трёх акустических ветвей со скоростью $v_0$
\[
    g(\omega)=\frac{3V\omega^2}{2\pi^2v_0^3}.
\]
В \rconcept{модели Дебая} спектр обрывается на частоте $\omega_D$,
при которой полное число мод равно $3N$:
\[
    3N=\int_0^{\omega_D}g(\omega)\,d\omega
      =\frac{V\omega_D^3}{2\pi^2v_0^3},\qquad
    \omega_D=v_0\left(\frac{6\pi^2N}V\right)^{1/3},\qquad
    \Theta=\hbar\omega_D.
\]
Тогда $g(\omega)=9N\omega^2/\omega_D^3$ и, с $z=\hbar\omega/T$,
\[
    E=\int_0^{\omega_D}\frac{\hbar\omega\,g(\omega)\,d\omega}
        {e^{\hbar\omega/T}-1}
     =\frac{9NT^4}{\Theta^3}\int_0^{\Theta/T}\frac{z^3\,dz}{e^z-1}.
\]

\subsubsection*{Высокие температуры, $T\gg\Theta$}
\[
    E\simeq\frac{9NT^4}{\Theta^3}\int_0^{\Theta/T}z^2\,dz=3NT,
    \qquad C=3N.
\]
Это \rconcept{закон Дюлонга--Пти}.

\subsubsection*{Низкие температуры, $T\ll\Theta$}
Верхний предел уходит в бесконечность,
$\int_0^\infty z^3dz/(e^z-1)=\pi^4/15$:
\[
    E=\frac{3\pi^4}{5}\frac{NT^4}{\Theta^3},\qquad
    C=\frac{12\pi^4}{5}N\left(\frac T\Theta\right)^3.
\]

\begin{rbox}[left=18pt,right=18pt,top=12pt,bottom=10pt]{[]}
\centering
\begin{tikzpicture}
\begin{axis}[
    width=0.78\linewidth,height=6cm,
    xmin=0,xmax=1.5,ymin=0,ymax=1.1,
    xlabel={$T/\Theta$,\quad $T/\hbar\omega_0$},
    ylabel={$C/C_\infty$},
    axis line style={rtext},tick style={rtext},
    every tick label/.style={text=rtext},
    label style={text=rtext},
    legend style={draw=none,fill=none,text=rtext,
        at={(0.97,0.05)},anchor=south east,font=\small},
    legend cell align=left]
\addplot[rconceptcolor,thick,smooth] coordinates {
    (0,0) (0.01,0.0210) (0.03,0.0631) (0.05,0.1058) (0.08,0.1727)
    (0.10,0.2205) (0.13,0.2981) (0.16,0.3792) (0.20,0.4825)
    (0.25,0.5921) (0.30,0.6772) (0.40,0.7900) (0.50,0.8554)
    (0.60,0.8953) (0.80,0.9385) (1.00,0.9598) (1.20,0.9718)
    (1.50,0.9818)};
\addlegendentry{цепочка, $C/N$}
\addplot[rc3,thick,smooth] coordinates {
    (0,0) (0.01,0.0001) (0.03,0.0021) (0.05,0.0097) (0.08,0.0397)
    (0.10,0.0758) (0.13,0.1524) (0.16,0.2445) (0.20,0.3686)
    (0.25,0.5031) (0.30,0.6077) (0.40,0.7459) (0.50,0.8254)
    (0.60,0.8738) (0.80,0.9260) (1.00,0.9517) (1.20,0.9661)
    (1.50,0.9781)};
\addlegendentry{Дебай 3D, $C/3N$}
\addplot[rtext,densely dotted,domain=0:1.5] {1};
\end{axis}
\end{tikzpicture}
\captionof{figure}{Теплоёмкость цепочки и кристалла Дебая.
При низких температурах $C\propto T$ и $C\propto T^3$ соответственно,
при высоких обе стремятся к классическому пределу.}
\end{rbox}

\remark[Двумерная решётка.]
На занятии отдельно не разбиралась. С $g_{2\mathrm D}\propto\omega$
из прошлого занятия тот же расчёт даёт $C\propto T^2$ при низких
температурах и $C=2N$ (две ветви в плоскости) при высоких.

\answer
\[
    C_{1\mathrm D}=
    \begin{cases}
        \dfrac{2\pi}3\dfrac{NT}{\hbar\omega_0}, & T\ll\hbar\omega_0,\\[8pt]
        N, & T\gg\hbar\omega_0,
    \end{cases}
    \qquad
    C_{3\mathrm D}=
    \begin{cases}
        \dfrac{12\pi^4}5N\left(\dfrac T\Theta\right)^3, & T\ll\Theta,\\[8pt]
        3N, & T\gg\Theta.
    \end{cases}
\]
\rresult{При низких температурах $C\propto T^d$} для линейного
акустического спектра.
\endgroup

\Needspace{16\baselineskip}
\task[3]{Нижняя фононная ветвь графена}
\label{task:fksv_graphene}

\condition
Определить вклад в теплоёмкость от нижней фононной ветви графена
$\omega=\beta k^2$. Обсудить температурную зависимость.

\solution

\subsection*{Плотность состояний}

Для листа площади $S$
\[
    g(\omega)=\frac{S}{(2\pi)^2}\int_0^\infty
       \delta(\omega-\beta k^2)\,2\pi k\,dk
     =\frac S{2\pi}\cdot\frac{k}{2\beta k}
     =\frac S{4\pi\beta}.
\]
Здесь применено $\delta(f(k))=\delta(k-k_0)/|f'(k_0)|$
с $f'(k)=-2\beta k$.
\rresult{Плотность состояний постоянна}, как у одномерной цепочки
с линейным спектром. Ветвь содержит $N$ мод, поэтому граничная частота
\[
    \frac{S\omega_D}{4\pi\beta}=N,\qquad
    \omega_D=\frac{4\pi\beta N}S.
\]

\subsection*{Энергия}

С $z=\hbar\omega/T$, $z_D=\hbar\omega_D/T$:
\[
    E=\int_0^{\omega_D}\frac{\hbar\omega}{e^{\hbar\omega/T}-1}
       \frac{S\,d\omega}{4\pi\beta}
     =\frac{S\hbar}{4\pi\beta}\left(\frac T\hbar\right)^2
       \int_0^{z_D}\frac{z\,dz}{e^z-1}.
\]

\subsubsection*{Низкие температуры, $T\ll\hbar\omega_D$}
$\int_0^\infty z\,dz/(e^z-1)=\Gamma(2)\zeta(2)=\pi^2/6$:
\[
    E=\frac{\pi^2}6\,\frac{NT^2}{\hbar\omega_D},\qquad
    C=\frac{\pi^2}3\,\frac{NT}{\hbar\omega_D}
     =\frac{\pi S}{12\hbar\beta}\,T.
\]

\subsubsection*{Высокие температуры, $T\gg\hbar\omega_D$}
$e^z-1\simeq z$:
\[
    E=\frac{S\hbar}{4\pi\beta}\left(\frac T\hbar\right)^2z_D=NT,
    \qquad C=N.
\]

\subsection*{Температурная зависимость}

Показатель степени определяется поведением $g(\omega)$ на малых
частотах: при $g\propto\omega^s$ получается $C\propto T^{s+1}$.
Квадратичный спектр в 2D даёт $s=0$ и \rresult{линейную по $T$
теплоёмкость} вместо $T^2$ для линейных ветвей в плоскости.
Поэтому при низких температурах вклад нижней ветви графена
преобладает.

\answer
\[
    C=
    \begin{cases}
        \dfrac{\pi^2}3\dfrac{NT}{\hbar\omega_D}
          =\dfrac{\pi S}{12\hbar\beta}T, & T\ll\hbar\omega_D,\\[8pt]
        N, & T\gg\hbar\omega_D,
    \end{cases}
    \qquad \omega_D=\frac{4\pi\beta N}S.
\]

\end{document}
